\documentclass{article}

\usepackage[utf8]{inputenc}
\usepackage[T2A]{fontenc}
\usepackage[russian]{babel}

\begin{document}
\small

\section*{Task 1.2}

\subsection*{1) $x=1,\ r_y=0,\ r_z=0$?}

\[
A.1 \to A.2 \to A.3 \to B.1 \to B.2 \to B.3 \to A.4 \to B.4 \to A.5 \to B.5
\]

оба потока читают \texttt{x = 0} в \texttt{y} и \texttt{z}, и каждый из них записывает \texttt{x = 1}.

\subsection*{2) $x=2,\ r_y=0,\ r_z=1$?}

\[
A.1 \to A.2 \to A.3 \to A.4 \to B.1 \to B.2 \to B.3 \to A.5 \to B.4 \to B.5
\]

B читает \texttt{x = 0} и входит в \texttt{if}. A записывает \texttt{x = 1}. После этого B читает \texttt{x = 1} и записывает \texttt{x = 2}.
При этом \texttt{y = 0}.

\subsection*{3) $x=1,\ r_y=0,\ r_z=1$?}

Нет, это невозможно.

\begin{enumerate}
\item Предпологаем, существование трассы такой, что в конце получаем $x=1,\ r_y=0,\ r_z=1$.
\item Так как имем $r_z=1$ в конце исполнение,то это значит, что произошло $B.4$ и было прочитано $x=1$.
\item $x=1$ могло записаться только если было событие $A.5$, значит оно было раньше чем $B.4$: $A.5 < B.4$.
\item Мы знаем что $B.4$ всегда исполняется перед $B.5$.
\item Так как $r_z=1$, в $B.5$ записывается $x=z+1=2$, то есть после $B.5$ значение $x$ равно 2.
\item Чтобы в конце у нас получилось $x=1$, надо чтобы была запись $x=1$ после $B.5$, то есть $B.5$ дожно быть раньше $A.5$.
\item Получилось противоречие: $A.5 < B.4 < B.5$ и одновременно $B.5 < A.5$.
\item Следовательно, получить $x=1,\ r_y=0,\ r_z=1$ невозможно ни при каких трассах (в interleaving модели).
\end{enumerate}

\end{document}